\documentclass[a4paper,10pt]{article}
\usepackage[utf8]{inputenc}
\usepackage[T1]{fontenc}
\usepackage[french]{babel}
\usepackage{enumitem}
% Puces pleines au premier niveau, vides au second : jamais de tiret, confondu avec un signe moins.
\frenchsetup{StandardItemLabels=true}
\setlist[itemize,1]{label=\textbullet}
\setlist[itemize,2]{label=$\circ$}
\usepackage[table]{xcolor}
\usepackage{geometry}
\usepackage{tabularx}
\usepackage{xltabular}
\usepackage{amssymb}
\usepackage{amsmath}
\usepackage{textcomp}
\usepackage{fontawesome5}
\usepackage{tikz}

\geometry{hmargin=0.8cm, vmargin=1.5cm}

% Commande 6 cases d'auto-évaluation sur 2 lignes
\newcommand{\caseseval}{%
    \centering
    $\square\;\square\;\square$ \\
    \vspace{2mm}
    $\square\;\square\;\square$
}
\renewcommand{\tabularxcolumn}[1]{m{#1}}

\begin{document}

\begin{center}
    \LARGE \textbf{\textcolor{orange!80!black}{\faBookOpen} \quad FICHE MÉMO}\\[0.5em]
    \large \textbf{3e - Chapitre 10 - Calcul littéral (partie 2)}
\end{center}

\vspace{0.2cm}

% arraystretch : 2.2 par défaut ; 1.8 si la fiche gagne à être plus dense.
% Le tableau est continu : les objectifs s'enchainent et la coupure de page est automatique.
\renewcommand{\arraystretch}{2.2}
\begin{center}
    \begin{xltabular}{\dimexpr\textwidth-5pt\relax}{|m{2cm}|X|X|X|}
        \hline
        \rowcolor{gray!10}
        \centering \scriptsize \textbf{\textcolor{red}{\faTimes}/\textcolor{green!60!black}{\faCheck} Auto-évaluation} &
        \centering \textbf{\textcolor{blue!70!black}{\faQuestionCircle} Question} &
        \centering \textbf{\textcolor{red!70!black}{\faBullseye} Réponse} &
        \centering \textbf{\textcolor{orange!90!black}{\faLightbulb} Stratégie / Conseil} \arraybackslash \\
        \hline
        \endhead

        \multicolumn{4}{|c|}{\cellcolor{blue!10}\parbox{\dimexpr\textwidth-20pt\relax}{\centering \textit{\textbf{Objectif : Je suis capable de développer et réduire avec la double distributivité}}}} \\
        \hline
        \caseseval & La formule de la double distributivité est : $(a+b)(c+d) = \dots$ & $ac + ad + bc + bd$ & \footnotesize \textit{Procède par ordre : distribue d'abord $a$, puis distribue $b$ ensuite.} \\ \hline
        \caseseval & Que sous-entend l'écriture simplifiée d'une expression comme $ac$ ? & La présence d'un signe $\times$ ($a \times c$). & \footnotesize \textit{Le signe de multiplication disparait devant une lettre ou une parenthèse.} \\ \hline
        \caseseval & \textit{Exemple :} Développe $(x+3)(x+5)$. & $x^2 + 8x + 15$ & \footnotesize \textit{Distribue dans l'ordre, puis regroupe les termes en $x$. As-tu bien obtenu quatre produits avant de réduire ?} \\ \hline
        \multicolumn{4}{|c|}{\cellcolor{blue!10}\parbox{\dimexpr\textwidth-20pt\relax}{\centering \textit{\textbf{Objectif : Je suis capable de développer avec l'identité remarquable $(a-b)(a+b)$}}}} \\
        \hline
        \caseseval & L'identité remarquable se développe ainsi : $(a-b)(a+b) = \dots$ & $a^2 - b^2$ & \footnotesize \textit{Identifie précisément qui joue le rôle de $a$ et de $b$ avant de calculer leurs carrés.} \\ \hline
        \caseseval & L'ordre des parenthèses change-t-il le résultat ? Que donne $(a+b)(a-b)$ ? & Non. Cela donne toujours $a^2 - b^2$. & \footnotesize \textit{La multiplication est commutative. L'essentiel est de repérer les mêmes termes avec des signes opposés.} \\ \hline
        \caseseval & \textit{Exemple :} Développe $(7x-4)(7x+4)$. & $49x^2 - 16$ & \footnotesize \textit{Attention au piège : c'est le terme entier qui est élevé au carré. $(7x)^2$ n'est pas $7x^2$.} \\ \hline
        \multicolumn{4}{|c|}{\cellcolor{blue!10}\parbox{\dimexpr\textwidth-20pt\relax}{\centering \textit{\textbf{Objectif : Je suis capable de factoriser avec un facteur commun}}}} \\
        \hline
        \caseseval & Comment appelle-t-on l'action de transformer une somme ou une différence en un produit ? & Factoriser. & \footnotesize \textit{C'est l'opération exactement inverse du développement.} \\ \hline
        \caseseval & La formule de factorisation simple s'écrit : $k \times a + k \times b = \dots$ & $k(a+b)$ & \footnotesize \textit{La lettre $k$ représente le facteur commun que tu vas isoler devant les parenthèses.} \\ \hline
        \caseseval & Par quelle étape commence-t-on pour factoriser ? & Décomposer les termes pour faire apparaitre un facteur commun. & \footnotesize \textit{Cherche dans les tables de multiplication (nombres) ou repère une lettre répétée. L'as-tu souligné visuellement ?} \\ \hline
        \caseseval & \textit{Exemple :} Factorise $35t^2 + 15z$. & $5(7t^2 + 3z)$ & \footnotesize \textit{Cherche le plus grand diviseur commun aux deux coefficients. Vérifie en redéveloppant mentalement si tu retombes sur l'expression de départ.} \\ \hline
        \multicolumn{4}{|c|}{\cellcolor{blue!10}\parbox{\dimexpr\textwidth-20pt\relax}{\centering \textit{\textbf{Objectif : Je suis capable de factoriser avec l'identité remarquable $a^2 - b^2$}}}} \\
        \hline
        \caseseval & La différence de deux carrés se transforme en produit selon la formule : $a^2 - b^2 = \dots$ & $(a-b)(a+b)$ & \footnotesize \textit{Ouvre systématiquement deux parenthèses : l'une avec une soustraction, l'autre avec une addition.} \\ \hline
        \caseseval & À quelle condition stricte (quel signe) peut-on appliquer cette formule ? & Il faut un signe ``moins'' (une différence) entre les deux carrés. & \footnotesize \textit{Avec un signe ``plus'', cette méthode ne s'applique pas. C'est un piège classique.} \\ \hline
        \caseseval & Quelle astuce appliquer pour bien identifier les rôles de $a$ et $b$ ? & Réécrire chaque terme sous la forme d'un carré $(\dots)^2$. & \footnotesize \textit{Demande-toi pour chaque terme : ``Quel nombre ou quelle lettre, multiplié par lui-même, donne ce résultat ?''} \\ \hline
        \caseseval & \textit{Exemple :} Factorise $x^2 - 25$. & $(x-5)(x+5)$ & \footnotesize \textit{Demande-toi de quel nombre $25$ est le carré avant d'ouvrir les parenthèses.} \\ \hline
        \multicolumn{4}{|c|}{\cellcolor{blue!10}\parbox{\dimexpr\textwidth-20pt\relax}{\centering \textit{\textbf{Objectif : Je suis capable de résoudre une équation produit nul}}}} \\
        \hline
        \caseseval & Énonce la règle d'un produit nul. & Si un produit de facteurs est nul, alors au moins l'un de ses facteurs est nul. & \footnotesize \textit{Cette phrase théorique est le point de départ : elle justifie la séparation de ton équation.} \\ \hline
        \caseseval & En langage mathématique : si $A \times B = 0$, alors \dots & $A = 0$ ou $B = 0$. & \footnotesize \textit{Le petit mot ``ou'' est crucial. Il sépare ton problème en deux mini-équations faciles.} \\ \hline
        \caseseval & Que faut-il impérativement vérifier avant d'appliquer cette méthode ? & Qu'on a bien un produit égal à ZÉRO. & \footnotesize \textit{Regarde le membre de droite : sans un $0$, la méthode du produit nul ne s'applique pas en l'état.} \\ \hline
        \caseseval & \textit{Exemple :} Résous $(x-7)(5x+5) = 0$. & Les solutions sont $7$ et $-1$. & \footnotesize \textit{As-tu bien isolé $x$ dans chaque mini-équation, l'une après l'autre ?} \\ \hline
        \multicolumn{4}{|c|}{\cellcolor{blue!10}\parbox{\dimexpr\textwidth-20pt\relax}{\centering \textit{\textbf{Objectif : Je suis capable de résoudre une équation du type $x^2 = a$}}}} \\
        \hline
        \caseseval & Pourquoi l'équation $x^2 = a$ n'a-t-elle aucune solution si $a < 0$ ? & Parce que le carré d'un nombre est toujours positif. & \footnotesize \textit{Appuie-toi sur la règle des signes : un moins par un moins donne toujours un plus.} \\ \hline
        \caseseval & Quelle est l'unique solution de l'équation si $a = 0$ ? & $0$ & \footnotesize \textit{Le seul nombre qui, multiplié par lui-même, donne zéro, c'est zéro.} \\ \hline
        \caseseval & Quelles sont les DEUX solutions distinctes de l'équation si $a > 0$ ? & $\sqrt{a}$ et $-\sqrt{a}$ & \footnotesize \textit{L'erreur la plus fréquente est d'oublier la solution négative. Penses-y systématiquement.} \\ \hline
        \caseseval & \textit{Exemple :} Résous $x^2 = 9$. & Les solutions sont $3$ et $-3$. & \footnotesize \textit{Les deux nombres conviennent : leur carré vaut $9$. As-tu pensé au négatif ?} \\ \hline

    \end{xltabular}
\end{center}
\end{document}
